\par
\specialsection{Обзор литературы}

Задача выбора следующей культуры в литературе рассматривается в нескольких близких, но не совпадающих постановках.
Наиболее часто встречаются работы по оптимизации севооборота, по crop recommendation на основе методов машинного
обучения и по системам поддержки принятия решений, в которых существенную роль играют интерпретация результата и форма
его представления пользователю \cite{fruchtfolge2021,cropml2024,agri-dss-visual,c3po2023}.

В первой группе работ основное внимание уделяется построению систем, ориентированных на \textbf{планирование и
оптимизацию} структуры посевов. В таких исследованиях задача формулируется как поиск наилучшей последовательности
культур или наилучшего распределения посевов с учетом множества ограничений. Примером является система
\emph{Fruchtfolge}, в которой выбор культур описывается как задача пространственно-явного оптимизационного планирования
с использованием математического программирования \cite{fruchtfolge2021}. Подобные решения близки к задаче полного
аграрного планирования, однако они отличаются от более узкой постановки, в которой по истории конкретного поля требуется
сформировать ранжированный список вероятных следующих культур.

Во второй группе работ задача ставится как \textbf{crop recommendation} на основе данных о поле, регионе, почве, климате
или истории наблюдений. Здесь акцент делается на построении модели, которая по набору признаков рекомендует наиболее
подходящую культуру или ранжирует несколько кандидатов \cite{cropml2024,xai-crop-recommendation2024}. Для данного
направления характерно использование методов машинного обучения как основного предиктивного ядра. При этом в современных
исследованиях все чаще подчеркивается, что для практического применения системы одной только точности недостаточно:
пользователь должен понимать, по каким причинам тот или иной вариант рекомендован, и насколько надежен результат
\cite{agri-dss-visual}, \cite{xai-crop-recommendation2024}.

Отдельное направление составляют работы, в которых используются \textbf{графовые и онтологические представления знаний}.
В подобных исследованиях knowledge graph выступает не столько как самостоятельная предиктивная модель, сколько как
средство описания сущностей предметной области, их связей, ограничений и правил интерпретации. Показательным примером
является проект \emph{C3PO}, в котором crop planning and production process ontology и knowledge graph используются как
основа knowledge-based decision support systems для сельского хозяйства \cite{c3po2023}. Для данной дипломной работы
этот класс решений важен потому, что показывает возможность применять графовое представление как интерпретационный слой,
дополняющий основное модельное ядро.

Существенным аспектом agricultural DSS является не только вычислительная часть, но и \textbf{форма представления
результата}. В обзоре визуализаций для agricultural decision support systems отмечается, что карты являются одним из
центральных интерфейсов таких систем, а также подчеркивается растущая значимость dashboard-подходов и визуализации
неопределенности \cite{agri-dss-visual}. Это особенно важно в задачах, где система формирует не единственный ответ, а
несколько кандидатов, требующих последующего анализа и сравнения.

Если рассматривать \textbf{прикладные существующие решения}, то на рынке присутствуют крупные цифровые платформы
управления полями и хозяйством. Например, John Deere Operations Center позиционируется как онлайн farm management
system, позволяющая планировать, контролировать и анализировать работу хозяйства, включая управление полями, техникой и
производственными операциями \cite{john-deere-operations-center}. Аналогично xarvio FIELD MANAGER ориентирован на
поддержку crop production decisions, используя spatial и satellite-based данные для анализа полевых зон, seeding, crop
nutrition и crop protection \cite{xarvio-field-manager}. Однако подобные платформы охватывают значительно более широкий
круг задач, чем рекомендация следующей культуры по истории поля, и потому не являются прямыми аналогами рассматриваемой
работы. Корректнее рассматривать их как более широкую прикладную среду, в рамках которой отдельный рекомендательный
модуль подобного типа потенциально может использоваться.

Таким образом, существующие исследования и прикладные решения показывают, что задача выбора культуры может
рассматриваться в разных масштабах: от полного оптимизационного планирования до узких рекомендательных и
интерпретационных модулей \cite{fruchtfolge2021,cropml2024,agri-dss-visual,c3po2023,xai-crop-recommendation2024}. На
этом фоне данная работа занимает промежуточное положение. Она не ставит целью построение полного оптимизатора
севооборота и не ограничивается только одновариантным прогнозом следующей культуры. Ее особенность состоит в сочетании
предиктивного ядра, Top-k рекомендательного режима, confidence-сигнала и graph-based explainability layer. В результате
работа ориентирована на построение \textbf{объяснимого рекомендательного прототипа}, который формирует shortlist
вероятных следующих культур и предоставляет пользователю интерпретируемый контекст для их анализа.
